\title{X-LoRA: Mixture of Low-Rank Adapter Experts, a Flexible Framework for Large Language Models with Applications in Protein Mechanics and Molecular Design}

\begin{abstract}
We introduce a mixture of expert approach to develop fine-tuned large language models using a deep, layer-wise token-level technique grounded in low-rank adaptation (LoRA). Beginning with a collection of pre-trained LoRA adapters, our gating mechanism leverages hidden states to dynamically blend adapted layers, enabling the resulting X-LoRA model to utilize diverse capabilities and generate novel deep layer-wise combinations for solving tasks. This design is inspired by biological concepts of universality and diversity, where neural network components are reused across different hierarchical levels. Consequently, the X-LoRA framework can be seamlessly applied to any existing large language model (LLM) without requiring changes to the base architecture. We create a specialized X-LoRA model equipped with scientific functionalities, including forward/inverse analysis tasks and improved reasoning skills, primarily targeting biomaterial analysis, protein mechanics, and design. The significance of this work lies in providing accessible, extensible, and adaptable models endowed with strong domain expertise and the ability to integrate knowledge across disciplines. Incorporating experts in biology, mathematics, reasoning, bio-inspired materials, mechanics and materials science, chemistry, protein biophysics, mechanics, and quantum-mechanics-based molecular properties, we perform a series of physics-oriented case studies. These include assessments of knowledge retrieval, protein mechanics forward/inverse problems, protein design, adversarial agentic modeling such as ontological knowledge graph creation, and molecular design. The model is not only capable of quantitatively predicting nanomechanical properties of proteins or quantum mechanical molecular features but also interprets the results and accurately forecasts probable mechanisms that underlie distinct molecular behaviors.
\end{abstract}

\section{Introduction}
Large language models (LLMs)~\cite{Vaswani2017AttentionNeedc,Touvron2023LlamaModels,OpenAI2023GPT-4Report,Chowdhery2022PaLM:Pathways,Jiang2023Mistral7B,Gunasekar2023TextbooksNeed} have become highly popular, including the creation of specialized models that excel in particular types of tasks, reasoning, or scientific fields~\cite{Bubeck2023SparksGPT-4,Buehler2023MechGPTModalities,Nejjar2023LLMsAnalysis,Buehler2023GenerativeDesign,Luu2023BioinspiredLLM:Materials,Luu2023GenerativeSolvents,Buehler2023MeLMProblemsb,Ge2023OpenAGI:Experts}. However, training such models can be expensive, especially when a broad range of abilities is required. Techniques like low-rank adapters (LoRA)\cite{hu2021lora} have been suggested as a more resource-efficient alternative, but these adaptations typically target more specific knowledge areas. The core idea behind LoRA is to introduce low-rank matrices added to the original full-size matrix, and to treat these low-rank matrices as the sole trainable elements of the model. Since only the adapter layers are updated during training, these models tend to retain the knowledge acquired during the base model’s pre-training while adapting the model for particular tasks, all while maintaining computational efficiency. Because training LoRA adapters is efficient, this approach allows for the development of a large number of specialized models.

Although LoRA adapters offer an efficient means to build specialized capabilities, an unresolved issue is how to combine multiple such systems into enhanced models that provide a unified and potentially superior set of functionalities. In fact, combining LLMs has become an intriguing topic of investigation~\cite{Kim2023SOLARUp-Scaling}. Other exploratory methods, including the use of spherically interpolated model parameters (SLERP) and related techniques~\cite{Arcee-ai/mergekit:Models.}, have demonstrated promising outcomes and the potential for combining models to attain novel capabilities. Models created through these mixing techniques are typically generated {\it a priori} following a specific algorithm, and while they have shown encouraging results on certain benchmarks, further study is required to systematically understand the construction of mixed models and the emergence of particular features.

We propose that dynamically mixing parameters could serve as a more universal method, as it enables the mixed model to explore novel combinations during inference.

Similar methodologies have employed mixture of experts (MoE) frameworks~\cite{Jacobs1991AdaptiveExperts,Hampshire1992TheRecognition,Jordan1994HierarchicalAlgorithm}, including the recently introduced Mixtral LLM~\cite{Jiang2024MixtralExperts}. Nevertheless, these approaches tend to be computationally intensive, requiring substantial resources even at inference time. In this work, we introduce a computationally efficient mixture-of-expert technique that utilizes a collection of distinct LoRA adapters, each of which can be trained independently with ease. Our method features a fine-grained and adaptable scheme, inspired by a biological model that reuses neural network components to build a hierarchy of interacting models, ranging from the foundational model to agent-based systems (Fig.~\ref{fig:hierarchical_concept}). To quantify its effectiveness, our study concentrates on applying this strategy to science-driven LLMs, specifically addressing protein mechanics challenges within the broader domain of bio-inspired materials analysis and design.

We develop an algorithm to dynamically combine adapters at the token-level state with fine granularity, enabling the mixing of adapters across individual layers. This approach, termed X-LoRA, provides access to a diverse array of capabilities that are essential for scientific applications of multimodal language modeling~\cite{Hu2023DeepScience,Buehler2023AMetamaterials,Buehler2022ProteinNetworks,}. 

\subsection{Fundamental concepts of X-LoRA}

We begin with a brief review of the concepts underlying the development of LoRA adapters. The core idea behind low-rank adaptation (LoRA)~\cite{hu2021lora} is based on the hypothesis that the weight updates lie in a low “intrinsic dimension.” This insight is leveraged by freezing the original weights $W_0 \in \mathbb{R}^{d \times k}$ and restricting the updates to a low-rank factorization

\begin{equation}
W_0 + \Delta W_0 = W_0 + BA 
\end{equation}

where $B \in \mathbb{R}^{d \times r}$ and $A \in \mathbb{R}^{r \times d}$ with rank $r \ll {\rm min}(d,k)$. During LoRA~\cite{hu2021lora} training, the pretrained weights $W_0$ remain fixed, and only the matrices $A$ and $B$ contain trainable parameters. The forward computation of a LoRA layer can be represented as:

\begin{equation}
    h = W_0x + \Delta W_0x = W_0x + BAx
\end{equation}

In practical applications, LoRA adapters include a fixed scalar weight $\alpha$, which is applied to the decomposed matrices, resulting in

\begin{equation}
    h = W_0x + BAx \cdot \alpha
\end{equation}

Expanding on this concept, X-LoRA utilizes a collection of LoRA adapters, each endowed with distinct functionalities, as demonstrated in prior science-oriented LLM research (e.g., \cite{Buehler2023GenerativeDesign}). Rather than statically combining or integrating models, we introduce a dynamic gating mechanism that adjusts each LoRA adapter individually at the granularity of tokens and layers, enabling fine-grained mixing deep within the model. More specifically, for a LoRA adapter $i$, the original scaling parameter $\alpha_i$ is multiplied by a factor $\lambda_i$ to yield a new scaled value $\alpha^*_i$:

\begin{equation}
    \alpha^*_i = \alpha_i \cdot \lambda_i
\end{equation}

The scaling factor $\lambda_i$ is predicted by an X-LoRA scaling head that leverages the model’s hidden states, forming the trainable part of the X-LoRA framework. Because the scaling head can utilize complex representations, it is efficiently implemented as a straightforward feed-forward neural network with one or a few layers (users can specify the architecture details of this network in X-LoRA’s implementation).

This scaling is computed separately for each token in the sequence, thus providing a high degree of granularity by scaling each LoRA adapter at specific LLM layers accordingly.

From here on, we refer to this method of selectively gating components of the adapters as scaling. Additional details on this technique are presented in the Materials and Methods section.

\subsection{Paper outline}

The structure of this paper is outlined as follows. Initially, we describe the methodology and training strategy, including the relevant datasets that contribute to the creation of an X-LoRA model with unique abilities in the physical sciences, with a particular emphasis on biomaterials such as proteins. Subsequently, we showcase a series of experiments applying the X-LoRA to diverse tasks, including question answering, conversational and agent-based modeling, protein design and analysis, among others. We provide a comprehensive examination of the scaling behaviors and validate the methodology by comparing it with molecular modeling and other physical data and techniques.

\section{Results and discussion}

The development of the X-LoRA model involves the following steps:

\begin{enumerate}
    \item Training the foundational base LLM (completed in prior work) \item Individually training multiple adapters to cultivate expertise in various fields (these correspond to distinct or overlapping domains of skills and knowledge) \item Training the combined X-LoRA model
\end{enumerate}

Our experiments begin with training a collection of LoRA adapters. We create nine adapters, each fine-tuned with specialized expertise, based on the Zephyr-7B-$\beta$ model~\cite{HuggingFaceH4/zephyr-7b-betaFace}, which itself is built upon the Mistral-7B model~\cite{Jiang2023Mistral7B}:

\begin{enumerate}
    \item Bioinspired materials \item Chain-of-thought (CoT) and reasoning \item Chemistry \item Mathematics \item Physics \item Biology \item Mechanics and materials \item Platypus: Logic and reasoning \item Protein mechanics tasks (including generative sequence-to-property and inverse capabilities)
\end{enumerate}

These specific domains of focus and expertise were selected because our ultimate goal is to develop a STEM-oriented LLM with an emphasis on materials science. Nevertheless, we hypothesize that our methodology is broadly applicable across various fields, particularly for generating modalities that facilitate connections between different knowledge domains~\cite{Buehler2023MechGPTModalities,Buehler2023MeLMProblemsb}.

Following the creation of the LoRA adapters, we proceed to train the X-LoRA model. This training utilizes a subset derived from the combined collection of training samples originally employed to train the individual LoRA adapters (further details are provided in Materials and Methods, including comprehensive descriptions of the dataset, implementation, and related aspects).

We treat the X-LoRA framework as an experimental platform to evaluate its effectiveness in multiple scenarios. Our experimental suite includes complex question answering, leveraging the generative model for protein design, performing analysis and reasoning, as well as devising more advanced interaction techniques through adversarial agentic modeling. During this evaluation, we compare the performance of the baseline model with that of the X-LoRA model and monitor how the scaling head engages various layers from the distinct experts during the generation or formulation of responses.

\subsection{Question answering and observed X-LoRA layer-wise scaling weights}

Initially, we investigate the differences between the X-LoRA model and the underlying foundation model. We pose the following question to both models:

\begin{quote}
\small\texttt{What material behavior (stiffening vs softening) causes supersonic crack speeds in mode II cracks? Think step by step.}
\end{quote}

The base model replies:

\begin{quote}
\small\texttt{Supersonic crack speeds, also known as Rayleigh waves, occur in mode II cracks due to the stiffening material behavior of the material near the crack tip. In mode II fracture, the loading is shear-dominated and the crack propagates perpendicular to the direction of applied stress. Here's a step-by-step explanation:

As the crack tip nears the material, the stresses in the surrounding area rise sharply. This creates a zone of intense stress concentration known as the stress singularity.

Within this stress singularity, the material undergoes elastic deformation which results in increased stiffness, or stiffening behavior. This stiffening occurs because the material adjacent to the crack tip becomes increasingly constrained as the crack moves closer [...]}
\end{quote}

Drawing on our own domain expertise~\cite{Buehler2008AtomisticFailure}, this response is not only inaccurate but also unnecessarily verbose. Conversely, the X-LoRA model replies with:

\begin{quote}
\small\texttt{The material behavior that causes supersonic crack speeds in mode II cracks is stiffening.}
\end{quote}

This response is accurate, concise, and straightforward~\cite{Buehler2008AtomisticFailure}.

The contrast between the base model and the X-LoRA model becomes even more evident in specialized protein analysis and design tasks. The given task is:

\begin{quote}
\small\texttt{CalculateForceHistory< M S K S N T Y R M L V L L E D D T K I N K E D E K F L K G K P G K M H E F V D E L I L P F N V D E L D E L N T W F D K F D A E I C I P Y V E I L K E S G M K > }
\end{quote}

The base foundation model responds with:

\begin{quote}
\small\texttt{I'm sorry but the input provided is not a valid mathematical expression. It appears to be a sequence of characters that do not form a calculable function. Please provide a valid mathematical expression for me to calculate.

However, I can tell you that the sequence of characters you provided is actually a series of function and variable names that may be part of a larger program or script in a specific programming language. Without context, it's impossible to determine the intended calculation.

If you provide more information, such as the programming language and the context in which this expression appears, I may be able to provide more insight into its meaning and possible calculations. }
\end{quote}

In contrast, the X-LoRA model produces the following response, successfully completing the design task by generating the force-deformation curve linked to this protein:

\begin{quote}
\small\texttt{[0.013,0.066,0.078,0.045,0.072,0.126,0.127,0.107,0.119, 
0.074,0.106,0.098,0.103,0.128,
0.077, 0.069,0.131,0.108,0.144,0.174, 0.166,0.183,0.185,0.217,0.175,0.230,0.187,0.224, 0.267,0.301,0.315]}
\end{quote}

Figure~\ref{fig:QA_weights_sample_0} presents outcomes from this question answering alongside the observed X-LoRA scaling weights. The two queries—one addressing mechanics and materials (Fig.~\ref{fig:QA_weights_sample_0}a) and the other focusing on protein mechanics (Fig.~\ref{fig:QA_weights_sample_0}b)—exhibit distinct activation patterns. Notably, a complex pattern of scaling values is evident in both scenarios, indicating that the X-LoRA model utilizes a heterogeneous combination of adapters across layers. Furthermore, the heatmap in Fig.~\ref{fig:QA_weights_sample_0}a reveals a wider distribution of activation across numerous experts and layers, including zones of high activation. Conversely, the heatmap in Fig.~\ref{fig:QA_weights_sample_0}b displays a more concentrated activation pattern, with fewer regions of elevated values. This observation suggests that the gating function on the right operates more selectively in activating experts.

In Fig.~\ref{fig:QA_weights_sample_0}, the areas resembling bright lines indicate paths of high activation, signifying that the scaling function selectively activates certain experts more strongly for particular layers. The observed patterns imply that the importance is not evenly distributed across all experts and layers; rather, it tends to be quite sparse and selective. Such sparsity aligns with findings in other mixtures-of-experts models, where only a subset of experts is engaged for a given input to specialize in different segments of the problem space.

By comparing these maps across various tasks, both the sparsity and the patterning may be crucial for better understanding which experts the model depends on for different input types or at distinct stages of neural network processing (the dynamic evolution of these maps during complex task solving will be addressed later). 

To deepen the analysis, we investigate a broader range of questions to assess how the X-LoRA model leverages different experts when presented with diverse prompts. Our focus is on whether, and if so how, the X-LoRA model integrates multiple adapters when questions span overlapping skills or domains. Fig.~\ref{fig:QA_weights} displays the outcomes from question answering along with the observed X-LoRA scaling weights, corresponding to the queries listed in Table~\ref{tab:table_qa}. In Table~\ref{tab:table_qa}, we present responses from both X-LoRA and the base model, Zephyr-7B-$\beta$. The distinctions between the two models are quite pronounced.

To highlight this visually, incorrect or dubious responses are indicated with a \redhighlight{red highlight}. The comparison demonstrates that X-LoRA achieves significantly higher accuracy alongside more succinct answers.

As predicted, we notice intricate mixing of adapters with frequent activation of multiple dominant LoRA experts. For example, Fig.~\ref{fig:QA_weights}d shows that a question regarding the mechanics of pine cones triggers strong activation of the mechanics/materials adapter in conjunction with the CoT/reasoning and bioinspired adapters. While X-LoRA selects the correct letter, the explanation is not entirely accurate since pine cones open in dry, rather than wet, conditions. However, a follow-up query probing the conditions for release produces the correct answer.

For a protein-centric question, Fig.~\ref{fig:QA_weights}f reveals adapter mixing in response to an inquiry about the mechanics of Bombyx mori silk, activating a blend of the bioinspired, mechanics/materials, and biology adapters. In tasks that are clearly specialized, such as protein design (Fig.~\ref{fig:QA_weights}a), we observe predominantly singular activation of a specific expert.

A more detailed examination of the heatmap of scalings shown in Fig.~\ref{fig:QA_weights}a reveals that, beyond the clearly prominent importance of the protein mechanics expert, there are several bands where particular experts receive significant emphasis across multiple layers. This pattern is especially evident for the mechanics/materials and reasoning experts. These bands are intermittent and exhibit a fluctuating pattern, which may indicate that these experts are particularly effective under certain conditions or for specific features. Notably, the reasoning expert (positioned immediately to the left of the protein mechanics expert) stands out with several layers where it exhibits the highest activation, particularly in the upper layers (approximately layers 100 to 140). We hypothesize that this could suggest a key role in the model's final decision-making or output generation, though further investigation is needed to confirm this. Additionally, there is a strong but more localized activation for the physics and biology experts at certain layers, potentially reflecting a specialization of these experts for features or representations in those areas.

Moreover, Fig.~\ref{fig:QA_token_history} illustrates a token-level history based on the aggregated scalings across all layers. This analysis presents histories for two cases: Fig.~\ref{fig:QA_token_history}a involves a task related to pine cone seed release, while Fig.~\ref{fig:QA_token_history}b shows the outcome of a sequence of protein mechanics tasks. The protein mechanics example consists of a series of prompt-answer exchanges: initially two property calculations followed by a generative task to design a novel protein towards the end. Although the utilization of experts remains generally consistent, there are noticeable shifts in the engagement of the LoRA experts throughout the process. Specifically, the protein mechanics expert becomes highly active during the final stages of generation when a new protein is being designed.

Table~\ref{tab:table_qa} presents a variety of use cases spanning multiple domains along with a comparison to the performance of the base model alone. We provide a more detailed discussion of several of these cases below.

For example, we present a question that combines elements of protein science and biology with logical reasoning, focusing on protein expression patterns. X-LoRA provides the correct answer. In contrast, the base model's response contains some errors and logical flaws when evaluating the expression of marker proteins in each cell type. The base model accurately recognizes that fibroblasts express Protein B and do not express Protein A due to the mutual exclusivity between Protein A and Protein B. However, it fails to explicitly mention that fibroblasts also express Protein C. Since there is no rule preventing the co-expression of Protein C with Protein B, and fibroblasts have no limitation regarding Protein C, it follows logically that they express Protein C as well. Regarding neurons, the base model begins correctly by indicating that neurons express Protein A and not Protein B, adhering to the mutual exclusivity rule. Yet, the subsequent discussion about neurons potentially expressing Protein C or both Protein C and Protein B is unclear and contains an error. Given that neurons express Protein A and not Protein B, the only remaining possibility consistent with the rules is that they express Protein C. There is no valid scenario in which neurons express Protein B, contrary to what was suggested initially. Therefore, neurons express Protein A and Protein C, but not Protein B. For epithelial cells, the base model correctly states that epithelial cells do not express Protein A but do express both Protein B and Protein C. Nevertheless, the explanation is convoluted and includes incorrect claims about the mutual exclusivity between Protein A and Protein C, which was not part of the original conditions. The straightforward reasoning is that epithelial cells lack Protein A expression but must express the other two proteins, Protein B and Protein C, since the initial instructions only specified mutual exclusivity between Protein A and Protein B.

X-LoRA, on the other hand, accurately determines the marker proteins expressed by each cell type based on the given observations and logical reasoning. For fibroblasts, X-LoRA correctly infers that fibroblasts do not express Protein A because the presence of Protein A excludes Protein B expression. Since fibroblasts do express Protein B, and considering that each cell type expresses a unique combination of three marker proteins (an incorrect assumption introduced, as the original statement did not specify the number of proteins each cell type expresses but only that they express a distinct combination), the conclusion that fibroblasts also express Protein C is valid according to the provided information (aside from the misunderstanding about the number of proteins). Regarding neurons, the model correctly identifies them as expressing Protein A and Protein C. The logic is valid, noting that Protein B cannot coexist with Protein A, and since neurons express Protein A along with one other marker protein (not Protein B), they must express Protein C. Lastly, for epithelial cells, the X-LoRA model correctly deduces that these cells express Protein B and Protein C. This is because epithelial cells lack Protein A expression, and by elimination, given that they express two marker proteins, these must be Protein B and Protein C.

Concerning the question on pine cone opening in humid versus dry conditions, the base model provides a confusing and factually incorrect response. X-LoRA delivers a clear and precise answer with well-explained reasoning. Similar examples include the question about the extraordinary mechanical properties of Bombyx mori silk, where the base model mistakenly claims that silk fibers conduct electricity. The base model also struggles with several other specialized domain questions, such as the role of hyperelasticity in maximum fracture speeds, where X-LoRA offers a concise and correct answer, but the base model fails to respond adequately. In a question about the strong adhesion of gecko feet, X-LoRA accurately highlights the significance of setae and nanoscale effects, whereas the base model wrongly attributes the process to special glue proteins. These and other instances, presented in the table and throughout the paper, provide compelling evidence of X-LoRA's superior performance compared to the base model.

Figure~\ref{fig:perf_comp_bioinspiredexam} presents a comparison of experiments evaluating knowledge recall using the bio-inspired knowledge test set introduced in \cite{Luu2023BioinspiredLLM:Materials}. It is evident that X-LoRA delivers the best results, despite being a considerably smaller model than either the BioinspiredLLM or Llama-BioLLM models (X-LoRA has 7B parameters compared to 13B in the others). The figure also includes comparisons with several other recently released LLMs, specifically Orca-13B and Llama-13b-chat.

Additionally, the plot compares X-LoRA against the base Zephyr-7B-$\beta$ model, demonstrating that X-LoRA exhibits substantially enhanced performance. These results suggest that excellent performance can be achieved with X-LoRA even when utilizing smaller base models.

Figure~\ref{fig:image_synth} illustrates an application in image synthesis, where X-LoRA is used to generate a prompt for other models, such as DALL-E 3 in this example.

\begin{quote}
\small\texttt{Create a prompt that I can use for image generation that accurately describes the microstructure of a hierarchical composite.

Explicitly describe the geometry.}
\end{quote}

The comparison shown in Figure~\ref{fig:image_synth}(b)-(c) clearly demonstrates that X-LoRA produces a more concise prompt, which in turn leads to a significantly more accurate image generated by DALL-E 3~\cite{DALLECard}. Importantly, the prompt created by Zephyr-7B-$\beta$ fails to precisely follow the given instruction and introduces additional design elements like fibers and nanoparticles, which were not mentioned in the original prompt focused specifically on the microstructure of a hierarchical composite. The prompt from X-LoRA is shorter and more focused, resulting in a much better output.

\subsection{Generative solutions to protein analysis}

We now turn to specific generative and analytical tasks related to proteins. These capabilities arise in the X-LoRA model through the protein mechanics adapter and are further strengthened by integrating the model's knowledge in biology, bio-inspired materials, reasoning, and logical deduction. This section concentrates on analyzing protein-related tasks independently and quantitatively assessing the performance of these abilities.

Fig.~\ref{fig:protein_force_history_prediction} presents findings from the protein task aimed at predicting force-deformation behaviors based on sequences. Fig.~\ref{fig:protein_force_energy_prediction} illustrates the overall accuracy in forecasting unfolding force and unfolding energy from the sequence, comparing actual values with predicted ones (with an observed $R_2=0.85$). In general, the model demonstrates excellent performance and, as will be discussed in the following section, it also excels at the inverse task of design and cycle consistent evaluation \cite{Lu2023GenerativePropertiesb}. However, unlike previous studies that employed highly specialized GPT-type models to address these tasks, our X-LoRA model merges these targeted capabilities with a variety of other deep specializations introduced through multiple fine-tuned adaptations.  

\subsection{Protein design and analysis}

We now broaden the scope of test cases to encompass more integrative applications spanning several capability areas. In particular, the experiments in this section highlight how the X-LoRA model leverages a distinct combination of skills and how agentic modeling can generate even more complex and comprehensive responses. We start the discussion with a protein design challenge, where the model is tasked with creating a novel protein that satisfies a specified target force-deformation profile. These force-deformation responses have been investigated in previous studies~\cite{Ni2023ForceGen:Model} and correspond to measurements of force versus deformation as the protein is stretched at both ends (representing a classic protein unfolding experiment~\cite{Lu1998UnfoldingSimulation.}).

Fig.~\ref{fig:design_protein} illustrates the application of the generative protein task. We design a protein exhibiting a targeted force-deformation response (training data and boundary conditions derived from the original molecular dynamics (MD) simulations of protein pulling, refer to~\cite{Ni2023ForceGen:Model}), then evaluate the predicted sequence to determine if it satisfies the desired properties (Fig.~\ref{fig:design_protein}a). This cycle-consistent evaluation is crucial to verify that the inverse design aligns with the forward prediction.

Fig.~\ref{fig:design_protein}b presents the predicted sequence: \texttt{MSKSNTYRMLVLLEDDTKINKEDEKFLKGKPGKMHEFVDELILPFNVDELDELNTWFDKFDAEICIPYVEILKESGMK}. Additionally, Fig.~\ref{fig:design_protein}c displays a folded structure of the protein model, obtained via AlphaFold 2~\cite{Jumper2021HighlyAlphaFoldb}. To investigate the relationship of the designed protein with existing sequences, Fig.~\ref{fig:design_protein}d shows a Basic Local Alignment Search Tool (BLAST) Tree~\cite{Altschul1990BasicTool} analysis. Notably, the designed protein shares similarities with hypothetical proteins explored in previous studies, which have yet to be structurally characterized.

Leveraging agentic modeling and multiple X-LoRA models as components within a multi-agent framework, we conduct a further examination of the predicted sequence using our X-LoRA model. This approach benefits from the model’s improved ability to comprehend protein mechanics and integrate this understanding with broader knowledge domains. Using the aforementioned predicted sequence as input, we request a more in-depth analysis and design recommendations for sequence variants. We initiate the task as follows:

\begin{quote}
\small\texttt{I have identified this amino acid sequence: MSKSNTYRMLVLLEDDTKINKEDEKFLKGKPGKMHEFVDEL ILPFNVDELDELNTWFDKFDAEICIPYVEILKESGMK. 

How can I use it? 

Provide logical steps, discuss specifically the engineering of the problem. }
\end{quote}

Through multiple agentic interactions employing adversarial strategies (see Materials and Methods for details), we encourage the model to explore more profound and intricate aspects of the topic. The X-LoRA model summarizes this dialogue as follows (for the complete conversation, refer to the Supplementary Material, SI Conv Example 1):

\begin{quote}
\small\texttt{The dialogue explores potential uses of a particular amino acid sequence in protein design. This sequence comprises both hydrophobic and hydrophilic residues, indicating it might form a stable conformation and engage with other molecules. To design a protein featuring this sequence, computational approaches can predict secondary and tertiary structures; however, challenges include the intricate nature of protein folding and the high computational demands. Enhancements can be achieved by creating more efficient algorithms and integrating novel experimental findings. Confirming the reliability of computational models is essential to guarantee their precision and practical relevance. Methods such as cross-validation, regularization, early stopping, ensemble techniques, and transfer learning can be employed to avoid overfitting and underfitting. Additionally, model interpretation and explanation aid in guiding practical uses.}
\end{quote}

We also requested the model to generate a summary of main points, yielding:

\begin{quote}
\small\texttt{\begin{itemize}
    \item The amino acid sequence MSKSNTYRMLVLLEDDTKINKEDEKFLKGKPGKMHEFVDELI LPFNVDELDELNTWFDKFDAEICIPYVEILKESGMK shows promise for protein engineering applications owing to its combination of hydrophobic and hydrophilic residues. Computational tools can predict the protein’s secondary and tertiary structures, though complexity in folding and computational expense pose limitations.
    \item Ensuring computational model validation is vital to maintain accuracy and practical utility. Strategies like cross-validation, regularization, early stopping, ensemble approaches, and transfer learning help mitigate overfitting and underfitting.
    \item Understanding and explaining models further supports their application in practice.
\end{itemize}
}
\end{quote}

Ultimately, the primary conclusion is:

\begin{quote}
\small\texttt{The key insight is that the amino acid sequence MSKSNTYRMLVLLEDDTKINKEDEKFLKGKPGKMHEFVDELILPFNVDELDELNTWFDKFDAEICIPYVEILKESGMK holds promise for protein engineering applications due to its combination of hydrophobic and hydrophilic residues. To design a protein containing this sequence, computational approaches can be employed to predict its secondary and tertiary structures; however, challenges include the complexity of protein folding and the computational resources required. It is essential to validate computational models to confirm their accuracy and practical relevance, using methods such as cross-validation, regularization, early stopping, ensemble techniques, and transfer learning to avoid both overfitting and underfitting. Additionally, interpreting and explaining the models can aid in guiding real-world applications.}
\end{quote}

We can delve further into specific strategies for sequence engineering. To this end, we request the model to provide more detailed design recommendations aimed at enhancing the protein's stability:

\begin{quote}
\small\texttt{What modifications can I make to the following protein sequence to improve its stability: MSKSNTYRMLVLLEDDTKINKEDEKFLKGKPGKMHEFVDELILPFNVDELDELNTWFDKFDAEICIPYVEILKESGMK?}
\end{quote}

The summary of the exchange is (complete dialogue available in Supplementary Material, SI Conv Example 2):

\begin{quote}
\small\texttt{The discussion centers on strategies to increase protein sequence stability through the introduction of hydrophobic amino acids and disulfide bonds. The biologist explains that hydrophobic residues typically enhance protein stability, and disulfide bonds further contribute to this effect. The engineer requests specific examples of hydrophobic amino acids and guidance on the number of disulfide bonds to introduce. The biologist cites examples of proteins successfully stabilized via these modifications, including insulin, bovine pancreatic trypsin inhibitor, and green fluorescent protein. The conversation also considers potential challenges or disadvantages of adding hydrophobic residues or disulfide bonds, such as altered protein function, increased structural complexity, and the risk of aggregation. The biologist recommends addressing these issues through careful amino acid selection, computational modeling, and empirical validation.}
\end{quote}

The list of main points:

\begin{quote}
\small\texttt{\begin{itemize}
\item To enhance the stability of a protein sequence, hydrophobic amino acids and disulfide bonds can be incorporated.

Hydrophobic amino acids generally support protein stability, while disulfide bonds also contribute to increasing stability.

Proteins such as insulin, bovine pancreatic trypsin inhibitor, and green fluorescent protein serve as examples where these modifications have been successfully applied to stabilize the proteins.
\item Possible drawbacks or challenges associated with introducing hydrophobic amino acids or disulfide bonds include alterations in protein function, added complexity, and a risk of aggregation.
\item These issues can be mitigated by careful amino acid selection, computational modeling, and experimental verification.
\end{itemize}
}
\end{quote}

In summary, the crucial point is:

\begin{quote}
\small\texttt{The most important conclusion from the discussion is that to increase the stability of the protein sequence MSKSNTYRMLVLLEDDTKINKEDEKFLKGKPGKMHEFVDELILPFNVDELDELNTWFDKFDAEICIPYVEILKESGMK, one can introduce hydrophobic amino acids and disulfide bonds. This directly addresses the initial question.}
\end{quote}

From the standpoint of protein engineering, these recommendations are logical as they are expected to yield more stable protein designs. Specific examples and strategies are outlined that researchers can further investigate.

In another instance, we examine three sequences. First, we request the model to evaluate the strength of each sequence, then to reason about the outcomes: Identify the strongest candidate and provide an explanation for why this protein is likely the strongest. Additionally, we ask the model to outline the steps necessary to produce it in the laboratory. The three sequences (manually modified from one in the test set) are:

\begin{quote}
\small\texttt{
\begin{itemize}
    \item LNTWFDKFDAEICIPYVEILKESGMK \item AITWFDKFDAEICIPYVEALKIAAMK \item AAAAAAAAKFDAAEICIIAAAAAAAAKIAAMK
\end{itemize}
}
\end{quote}

This scenario evaluates whether the model can integrate various abilities and adapt dynamically during the conversation—that is, to flexibly shift between expertise such as protein mechanics, logical analysis, hypothesis generation, and more.

The dialogue is summarized in Figure~\ref{fig:conversation_evolve}, accompanied by an illustration of how the scalings evolve throughout the conversation (the analysis presents both the detailed layer-wise adapter scaling and a cumulative analysis to provide an estimate of the effective utilization of the different experts). The findings clearly indicate that the protein task is initially highly relevant, but gradually shifts towards a more dominant representation of the bioinspired adapter and the biology adapter. By the conclusion of the conversation, the biology adapter emerges as the most dominant.

To evaluate the predictions and assess the quality of the responses, the three proteins involved in this task are shown in Figure~\ref{fig:protein_visuals} (proteins folded using Alpha Fold 2~\cite{Jumper2021HighlyAlphaFoldb}, via ColabFold \cite{Mirdita2022ColabFold:All}). We observe that the most stable structure in the center exhibits the most organized secondary structure, consistent with the idea that it has the highest unfolding force, as predicted by the X-LoRA model in Figure~\ref{fig:conversation_evolve}. The other two proteins are less structured; the first is mainly helical with some turns, while the last is partially unstructured and contains a bend in its geometry. These characteristics account for the lowest unfolding force predicted by the model.

By examining the visual depictions of the folded structures, we verify that the most stable structure in the center exhibits the highest level of secondary organization, consistent with the idea that it produces the greatest unfolding force as predicted by the model. We observe that the other two models show less organization: the first is primarily helical with some turns, while the last is partially unstructured and displays a kink in its geometry. We suggest that these characteristics probably account for the lowest unfolding force observed, aligning with the predictions of the X-LoRA model.

In a second case study, we turn our attention to the energetic characteristics of proteins, specifically analyzing the unfolding energy required as the protein is extended from its initial conformation to a fully unfolded state~\cite{Ni2023ForceGen:Model}. We again evaluate three sequences (the first two are identical to those mentioned above, whereas the third is a longer construct combining the second sequence at its termini and the first sequence in the middle):

\begin{quote}
\small\texttt{\begin{itemize}
    \item LNTWFDKFDAEICIPYVEILKESGMK \item AITWFDKFDAEICIPYVEALKIAAMK \item AITWFDKFDAEICIPYVEALKIAAMKLNTWFDKFDAEICIPYVEILKESGMKAITWFDKFDAEICIPYVEALKIAAMK
\end{itemize}
}
\end{quote}

Initially, we ask the model to calculate the unfolding energy for each sequence, and subsequently to interpret the findings: given the amino acid composition, we request an explanation for why this particular protein is likely the most stable. Finally, we inquire about the probable functions that the most stable protein might possess.

Figure~\ref{fig:MJB_20} presents a transcript of a dialogue that engages multiple experts and synthesizes knowledge from these various areas, demonstrated here in the context of unfolding energy for three proteins. By the conclusion of the conversation, the CoT/reasoning adapter emerges as the dominant component, as the X-LoRA model addresses a query that involves reasoning about the results to predict probable protein structural features along with potential protein functions. In particular, the mechanistic inquiry into the basis of protein stability highlights the formation of hydrophobic interactions among nonpolar amino acids and hydrogen bonds among polar amino acids. The model forecasts that several hydrophobic residues (I, F, W, Y) and hydrogen-bonding residues (D, E, K) play significant roles in protein stability.

Additionally, the model anticipates that cysteine (C) residues within the sequence would result in disulfide bond formation, further reinforcing protein structure stability. Notably, the model accurately identifies crucial structural features that explain the high stability, as corroborated by structural analyses depicted in Fig.~\ref{fig:MJB_21}. This analysis confirms the presence of these features and clarifies why the other two proteins, which consist of shorter segments with varying helical domain stability, exhibit considerably lower unfolding energies.

\subsection{Adversarial agentic modeling to connect distinct scholarly disciplines and knowledge yielding ontological knowledge graph generation}

Given that our X-LoRA model possesses extensive knowledge spanning multiple fields, we can leverage it to explore connections among diverse concepts, knowledge repositories, and expert domains. In one study, we pose two queries to the model, each crafted similarly but emphasizing different perspectives. Subsequently, we generate triplets to construct an ontological knowledge graph, thereby refining the answers into more structured outputs~\cite{Giesa2012CategoryDesign,Spivak2011ReoccurringAnalogies}. The resultant graph delivers a cohesive interpretation of the derived insights and depicts relationships between concepts in an interpretable, mechanistic fashion.

We employ two adversarial agents to address the question. Agent 1 is labeled as a "critical physicist," while Agent 2 assumes the role of a "philosopher." Agent 1 initiates the question and is responsible for probing and formulating follow-up questions in response to each reply from Agent 2. As the dialogue progresses, it delves more deeply into the topic and reveals various facets. The query is deliberately selected to span multiple domains of disciplinary knowledge, aiming to evaluate the model's capacity to synthesize ideas and concepts.

The query reads:

\begin{quote}
\small\texttt{Consider an abstract representation of a classical musical composition as a graph of interacting notes, chords, melodies, rhythms and so on.

What happens if we would apply pressure, up to a point of failure, on this abstract concept of music?

Provide logical steps, mathematical reasoning, and discuss specifically the physics of the problem.}
\end{quote}

The complete conversation is included in the Supplementary Material (see, SI Conv Example 3). For conciseness, we present only the summary and main conclusion here, capturing the core elements:

\begin{quote}
\small\texttt{The dialogue explores the effect of applying pressure to an abstract representation of a classical musical composition, treating it as a physical system composed of interacting notes, chords, melodies, and rhythms. This pressure can be quantified in terms of stress, with the critical failure point identified when the stress attains a threshold value. To validate the mathematical model's predictions, a physical analog of the abstract musical composition may be constructed, allowing experimental results to be compared against theoretical expectations. The proposed mathematical model has limitations, including a simplified depiction of interactions, static assumptions, uniform distribution, and a linear correlation between stress and pressure. Enhancing the model's fidelity and applicability would require integrating more detailed and realistic representations, dynamic considerations, non-uniform distributions, and nonlinear components.

[...]

The key insight is that applying pressure to an abstract classical musical composition can be conceptualized as modifying the force fields among the interacting notes, chords, melodies, and rhythms. As pressure increases, the internal stress of the system rises accordingly, culminating in failure and collapse at a critical threshold. This explanation provides a physical rationale for how pressure impacts the abstract musical concept and how such effects can be expressed quantitatively via stress.}
\end{quote}

Although the generated text offers valuable perspectives, constructing an ontological representation of the data can facilitate interpretability and analysis~\cite{Buehler2023MechGPTModalities}. For graph generation, we concatenate the entire dialogue between the two agents and build a knowledge graph. An example of such a graph is depicted in Figure~\ref{fig:graph}. The resulting graph is partitioned into communities using the Girvan-Newman algorithm~\cite{Girvan2002CommunityNetworks}, yielding a total of 12 communities. The figure displays both the overall graph (Figure~\ref{fig:graph}a) and zoomed-in views with node labels (Figure~\ref{fig:graph}b-c).

\subsection{Creation of X-LoRA-Gemma Integrating Protein, Chemical, Bio-inspired, and Materials Mechanics Functionalities}
\label{gemma-section}

To demonstrate that the proposed methodology is effective across different base models with varying architectures, we developed another X-LoRA model, this time built upon the Gemma-7B-it model~\cite{Google/gemma-7b-itFace}. This version incorporates four distinct LoRA adapters, defined as follows:

\begin{enumerate}
    \item Bioinspired materials \item Mechanics and materials \item Protein mechanics tasks (including generative sequence-to-property and inverse design capabilities) \item Quantum-mechanics based molecular properties from QM9 (featuring generative SMILES-to-property and inverse design functionalities, with Table~\ref{tab:table_QM9} providing definitions and summaries of all properties calculated in this dataset)~\cite{Ruddigkeit2012EnumerationGDB-17,Ramakrishnan2015ElectronicSpace}
\end{enumerate}

The X-LoRA-Gemma model is developed similarly to the original X-LoRA but is designed to possess a unique set of abilities, notably incorporating the prediction of 12 quantum mechanical properties and the capacity to design molecules that satisfy these 12 quantum mechanical criteria. This newly introduced adapter is trained concurrently with the other adapters mentioned above, all based on the same datasets utilized in our previous experiments. It is important to note that while the main focus of this paper is on the Zephyr-based X-LoRA model, in this section we distinguish the models by referring to the original as X-LoRA-Zephyr and the Gemma-based variant as X-LoRA-Gemma.

Figure~\ref{fig:protein_force_energy_prediction} illustrates the overall performance of X-LoRA, showing predictions of unfolding force and unfolding energy from sequences ($R_2 = 0.87$), alongside its effectiveness in predicting quantum mechanical properties of molecules derived from the QM9 dataset ($R_2 = 0.96$). Both tasks exhibit outstanding performance, with results for the protein design task slightly surpassing those of the original X-LoRA model.

Next, we demonstrate the functionalities of X-LoRA-Gemma through a straightforward multi-agent design experiment. We select a small organic molecule, \texttt{CCCC12CC1C(=O)N2} (drawn from the QM9 test set). This molecule is a bicyclic compound incorporating a ketone (=O) and an amine (N) within the ring structure; refer to Figure~\ref{fig:QM9_design}(a) for its visualization. Initially, we prompt the X-LoRA-Gemma model with:

\begin{quote}
\small\texttt{Here is a molecule in SMILES code: CCCC12CC1C(=O)N2

What modifications should I make to enhance this molecule’s suitability for polymer design?}
\end{quote}

The answer is:

\begin{quote}
\small\texttt{To enhance the usefulness of this molecule for polymer design, you might:
\begin{enumerate}
\item Add a reactive functional group, like an amine or hydroxyl group, to enable additional modifications and crosslinking.
\item Incorporate a longer, more flexible carbon backbone to improve the molecule’s capacity to form polymers with diverse structures and characteristics.
\item Attach a side chain containing a functional group, such as a carboxyl group, to facilitate the formation of covalent bonds with other molecules and polymers.
\item Embed a cyclic structure into the molecule, for example, a lactam ring, to produce a polymer that is more rigid and stable.
\item Introduce a variety of side chains to generate a broad range of properties and functionalities within the polymer.
\item Utilize a varying number of carbons in the backbone to enable the creation of polymers with different molecular weights and attributes.
\item Add a crosslinking site within the molecule to yield a more stable and less reactive polymer.
\item Incorporate a repeating unit structure to design a polymer with a defined molecular weight and properties.
\item Combine several of the above approaches to develop a versatile and practical polymer suitable for multiple applications.
\end{enumerate}}
\end{quote}

Next, we concentrate on the initial suggestion.

\begin{quote}
\small\texttt{How can I enhance the overall reactivity of a molecule? I intend to use this molecule to synthesize a polymer material.

The molecule in question, represented in SMILES notation, is: CCCC12CC1C(=O)N2

I have the option to adjust the following properties:

Dipole moment, $\mu$ Polarizability, $\alpha$ Highest Occupied Molecular Orbital energy, HOMO Lowest Unoccupied Molecular Orbital energy, LUMO HOMO-LUMO gap: The energy difference between HOMO and LUMO, gap Electronic spatial extent Zero-point vibrational energy Heat capacity at constant volume Internal energy at 0 K Internal energy at 298.15 K Enthalpy at 298.15 K Free energy at 298.15 K What are the three most important properties I should modify, and should I increase or decrease each? }
\end{quote}

The feedback from X-LoRA-Gemma is as follows:

\begin{quote}
\small\texttt{\begin{itemize}
    \item  Reduction of HOMO-LUMO gap \item  Increase in dipole moment \item  Enhancement of polarizability
\end{itemize}
Adjusting these properties can boost the overall reactivity of the molecule, rendering it more appropriate for constructing polymer materials. }
\end{quote}

We leverage both the generative and predictive functionalities of X-LoRA-Gemma to tailor these molecular property changes. 

For the optimization, we implement a two-agent framework comprising two separate instances of X-LoRA-Gemma. One agent performs the generative task \texttt{GenerateMolecularProperties<...>} to create candidate molecules. Subsequently, the second agent calculates molecular properties via \texttt{CalculateMolecularProperties<...>} and evaluates the discrepancy between the desired and predicted properties. This cycle continues until either the predetermined maximum iteration count is reached or the mean-squared error falls below a set threshold.  

The outcomes of this procedure are presented in Figure~\ref{fig:QM9_design}. Panel (a) illustrates the target molecular design alongside its deviation from the original molecule’s properties. Panel (b) displays the top 20 molecular designs produced by the algorithm, while panel (c) shows the distribution of the mean-squared error, and panel (d) plots the mean-squared error progression over the generated candidates. The optimal design identified is \texttt{CC(C)C1=CC(=O)NC=C1}, which closely aligns with the target as depicted in panel (e). Lastly, panel (f) provides an optimized 3D molecular representation obtained via UFF~\cite{Rappe1992UFFSimulations} for structural optimization. 

To evaluate the molecule’s reactivity, we also calculate the Gasteiger charges~\cite{Gasteiger1980IterativeCharges} (visualized in Figure~\ref{fig:QM9_design}(f)). The molecule contains an aromatic ring with a ketone (C=O) group, which is more electrophilic relative to typical carbon-carbon or carbon-hydrogen bonds, increasing its susceptibility to nucleophilic attack. The presence of nitrogen within the aromatic ring and the ketone group forms a pyridone ring structure. This ring exhibits higher reactivity than a simple aromatic ring due to the heteroatom (nitrogen) and ketone functionalities introducing reactive sites, such as the carbonyl carbon available for nucleophilic addition or positions on the ring amenable to electrophilic substitution. The nitrogen atom, integrated into the heteroaromatic system, influences the molecule’s aromaticity and reactivity by modifying electron density distribution across the ring. This effect can alter how the molecule interacts with other chemical species, making the nitrogen-containing ring a locus for potential chemical reactions, particularly involving electrophilic attack facilitated by the electron-donating characteristics of nitrogen.

The carbon atom within the ring, adjacent to both an oxygen and a nitrogen atom, carries a partial positive charge (+0.25), indicating it is somewhat electron-deficient. This occurs because oxygen, being more electronegative, draws electron density toward itself, resulting in the carbon having fewer electrons. This effect is partially offset by the neighboring nitrogen atom, which is less electronegative than oxygen but more so than carbon. Nevertheless, the presence of the partial positive charge implies that this carbon is a probable site for nucleophilic attack, where an electron-rich nucleophile would be attracted to this electron-poor carbon.

The nitrogen atom exhibits a partial negative charge (-0.33), suggesting it is relatively electron-rich compared to its standard state. This electron richness arises from its lone electron pair and the influence of the aromatic ring's electron system. In aromatic heterocycles, the nitrogen lone pair can participate in the delocalized $\pi$-electron network, though the extent of this involvement depends on the ring's configuration and the electronegativity of neighboring atoms. The partial negative charge indicates that nitrogen is more nucleophilic, meaning it tends to donate electrons more readily, either by forming bonds with electrophiles or through protonation. Additional quantum mechanical studies are necessary to further explore this behavior.

The carbon carrying a partial positive charge represents a reactive center for nucleophilic attack. Various molecules may interact with this site via mechanisms where nucleophiles target carbons deficient in electrons.

The hydrogen atom bonded to the nitrogen, with a partial positive charge of 0.17, is an excellent candidate for forming hydrogen bonds with other molecules or atoms. Hydrogen bonding is a type of dipole-dipole interaction that occurs when a hydrogen atom attached to a highly electronegative atom (such as nitrogen, oxygen, or fluorine) interacts with another electronegative atom possessing a lone pair of electrons. Consequently, this molecule can potentially act as a hydrogen bond donor, which plays a significant role in determining its solubility in water and other polar solvents, as well as its interactions within biological systems.

The ketone group positioned next to the nitrogen-containing aromatic ring imparts distinctive reactivity that can be harnessed in the synthesis of polymers~\cite{McQuarrieSimonPhysicalChemistry,CareySundbergAdvancedOrganicChemistry,Varghese2022BeyondApplications,Varghese2022BeyondApplications}. This compound is classified as a lactam due to the cyclic amide formed between the ketone and the nitrogen within the ring. Lactams play a crucial role in polymer chemistry, especially in creating polyamides, a category of polymers widely utilized in material applications. The ketone moiety is capable of undergoing several chemical reactions beneficial to polymer chemistry, including nucleophilic addition. Such reactions can be employed to lengthen the polymer chain or introduce side groups that alter the polymer’s characteristics. The nitrogen atom in the aromatic ring notably influences the polymer’s electronic attributes and can engage in hydrogen bonding, which may affect the melting point, solubility, and mechanical behavior of the polymer. Furthermore, N-heteroaromatic compounds are recognized for their capacity to coordinate with metal ions, potentially enabling uses in catalysis or in materials exhibiting electronic or photonic functionalities. Additionally, this molecule may undergo polymerization via the amide (lactam) group; polyamides formed in this manner are renowned for their durability, elasticity, and resistance to wear and chemicals, and find applications across diverse domains ranging from textiles (such as nylon) to automotive components~\cite{Varghese2022BeyondApplications,Varghese2022BeyondApplications}.

Alkyl substituents can enhance the hydrophobic nature of the molecule, leading us to expect that polymers synthesized from monomers containing such alkyl groups will exhibit increased hydrophobicity. This alteration may influence their solubility across various solvents as well as their interactions with water. Such properties might be beneficial for uses that demand water resistance or particular solubility profiles. Additionally, alkyl groups can impact the mechanical characteristics of the polymer. For example, they may function as plasticizers within the polymer matrix, thereby improving the polymer's flexibility. This effect could result in materials that are more flexible or possess enhanced impact resistance, contingent upon the alkyl groups' length and branching.

\subsubsection{Additional benchmarking} We present further benchmarking results for the X-LoRA-Gemma model, illustrated in Figure~\ref{fig:perf_MMLU}, using a selected subset of the MMLU benchmark that concentrates on chemistry, physics, and biology (domains closely related to those our model was trained on)~\cite{Hendrycks2020MeasuringUnderstanding}. Our findings indicate that both X-LoRA models described here outperform their respective base models. Specifically, the X-LoRA models show higher accuracy compared to the base versions (Zephyr model: 54.6\% for the base model versus 56.6\% for the X-LoRA model; Gemma model: 40.6\% for the base model and 52.4\% for the X-LoRA-Gemma model). The highest overall accuracy is attained by the X-LoRA-Zephyr model, although both models deliver comparable performance.

\section{Conclusion}

Multimodal LLMs have numerous current and prospective applications in scientific fields, alongside a strong demand for highly specialized, domain-specific expertise. Beyond general language tasks, performing calculation or generative tasks—as exemplified here with protein mechanics—is essential~\cite{Luu2023BioinspiredLLM:Materials,Buehler2023MechGPTModalities,Buehler2023MeLMProblemsb}. X-LoRA, which builds upon open-source language models, allows for the dynamic blending of expertise and knowledge from various domains, as illustrated through several examples including Figure~\ref{fig:conversation_evolve}. This investigation revealed how the X-LoRA model autonomously and adaptively modifies the use of adapters and how it combines different adapters in a complex manner across the LoRA layers. The intricate mixing of LoRA layers appears consistently in all cases examined here, such as in Figures~\ref{fig:QA_weights_sample_0} and \ref{fig:QA_weights}. This phenomenon also presents an intriguing challenge to better comprehend the underlying mechanisms from a physics standpoint, viewing an LLM as a large-scale graph-forming system that leverages complex graph structures. Consequently, the blending of adaptation experts at deep layers, as implemented here, changes these mechanisms to facilitate the development of effective solutions.

The proposed method enhances the traditional inference approach of a single forward pass by incorporating two forward passes. This trains the model to first ‘consider’ the question and how it might reconfigure itself before generating a response. This approach represents a straightforward implementation of ‘self-awareness,’ enabling the model to adjust its own architecture to optimally address a task. As a result, even with a relatively modest parameter size of 7 billion, the model can reason across diverse scientific disciplines (including biological materials, mathematics, physics, chemistry, logic, mechanics, etc.), greatly improving its ability to produce novel solutions. Additionally, this framework could inspire alternative model designs, such as approaches that extend the reconfiguration strategy to sections of the model or to entire non-adapted systems. In such cases, the scaling head would dynamically reconfigure portions of a trained model or combine multiple models. As demonstrated by X-LoRA, these strategies can serve as powerful tools for broadening or integrating the functionalities of different models.

We conducted multiple performance evaluations, including a comparison with significantly larger models (Figure~\ref{fig:perf_comp_bioinspiredexam}), in which X-LoRA distinctly outperforms them. The comprehensive comparison presented in Table~\ref{tab:table_qa} demonstrated that X-LoRA achieves substantially better results across a range of tasks and application fields than the base model. Additional comparisons involved quantitative evaluations of numerical prediction tasks, such as protein mechanics and the quantum-mechanical properties of molecules, where X-LoRA exhibited high accuracy with $R_2$ values ranging from 0.85 to 0.96. For reference, these results surpassed other published findings, such as the MeLM model~\cite{Buehler2023MeLMProblemsb} which reported an $R_2$ of 0.78 for predicting protein unfolding force. The ability to accurately address multiple prediction challenges, alongside inverse design tasks and complex reasoning skills, was highlighted in the molecular design example (Section~\ref{gemma-section}). This example also demonstrated that X-LoRA can be effectively implemented using different base model architectures (Gemma versus Zephyr/Mistral). Future research could further extend this by developing X-LoRA models for larger base models.

A limitation of the proposed method is the necessity for two forward passes, potentially increasing computational overhead: one to compute hidden states that serve as inputs to the X-LoRA scaling head, and another with $\Lambda$ applied to the adapters. Nonetheless, the initial pass is utilized for the X-LoRA scaling head and thereby leverages the fundamental capabilities of the LLM, unlike approaches that require training a larger scaling head to learn the knowledge independently. To effectively deploy this approach in model-serving environments, separate key-value caches must be maintained for both the scaling and forward passes. To enhance inference speed, we developed \texttt{Mistral.rs}, a LLM serving platform implemented in Rust~\cite{matsakis2014rust} that supports X-LoRA and incorporates several optimization techniques. Despite this, a key benefit of our approach is its compatibility with any model without the need to restructure internal logic. We consider this a major advantage, particularly given its applicability to the extensive resources available in the Hugging Face ecosystem.

A key benefit of the proposed method is that it offers an easy implementation path for any pre-existing LLM (for example, since the X-LoRA code is publicly available, it can be readily applied to any model within the Hugging Face framework). Another benefit lies in how the scaling head leverages the inherent capabilities of the employed LLM, learning intricate techniques for optimally combining adapters to produce specific responses. When training X-LoRA, it is ideal to use complex examples of question-answer pairs or dialogues that enable the model to discover the most effective combinatorial approach. We conducted various analyses related to protein design and mechanics. For example, Figure~\ref{fig:MJB_20} and Fig.~\ref{fig:MJB_21} illustrate a stability evaluation of multiple protein sequences, including a comparative assessment, reasoning based on predicted outcomes, and a mechanistic interpretation that grounds the predicted behavior in fundamental chemical principles. As demonstrated in Fig.~\ref{fig:MJB_21}, this was directly validated through structural examination of the protein’s folded 3D conformation.

Regarding future directions, further research could investigate a broader array of adapter experts, expanding upon the initial expertise present in our LoRA collection. The X-LoRA model outperforms both the base model and a base model equipped with a single adapter, delivering responses to complex queries by harnessing specialized abilities contributed by the adapter ensemble. Moreover, there appear to be synergistic interactions among the adapters, as suggested by the intricate mixing of scaling weights across layers. This phenomenon warrants additional study and could be compared to techniques like SLERP, which blend models using alternative approaches. While we trained the X-LoRA scaling head using a subset of the data—several hundred samples from each original training set utilized to develop the adapters—more targeted construction of such training datasets could be devised. For instance, we observed that multiple-choice questions tend to activate reasoning-focused adapters, given their training on related datasets. Supplementary prompting strategies are necessary to guide the scaling head in comprehending specific domains essential for answering particular questions, potentially through designated system messages or similar techniques. Overall, crafting suitable training datasets remains a promising area for further investigation, including specialized methods to facilitate effective layer-wise scaling mixing. Our experiments tracking these mechanisms over extended conversations reveal that the model can dynamically adjust scaling to optimally address distinct tasks, which is an encouraging result.

We identified intriguing patterns in how experts are activated across layers, with notable insights detailed in the paper. This analysis has the potential for further expansion and could provide valuable understanding of the benefits and mechanisms behind mixing model components or sets of layers. Such exploration would be particularly compelling when the X-LoRA method is applied beyond protein mechanics and protein design, as demonstrated here. We defer this to future research.

Another significant direction involves deeper investigation into agentic modeling, as illustrated here by employing two X-LoRA models interacting adversarially. Subsequent studies could incorporate more than two models to enrich interactions and introduce additional functionalities, thereby advancing models into uncharted generative domains. This approach may also facilitate the development of methods that combine physics-based modeling or other validation procedures, utilizing code-writing/executing agents or agents employing function calling and other processing techniques~\cite{Ghafarollahi2024ProtAgents:Learning}.

\section{Materials and Methods}

The X-LoRA codebase and examples are accessible at \url{https://github.com/EricLBuehler/xlora}. Model weights and datasets related to the X-LoRA discussed herein, along with supplementary examples, can be found at \url{https://huggingface.co/lamm-mit/x-lora}. The X-LoRA-Gemma model is available at \url{https://huggingface.co/lamm-mit/x-lora-gemma-7b}. The \texttt{Mistral.rs} inference engine is provided at \url{https://github.com/EricLBuehler/mistral.rs}.

\subsection{Mathematical formulation of X-LoRA} We define the $\mathbf{scalings}$ as a matrix that contains token- and layer-level coefficients for LoRA adapters. The X-LoRA $\mathbf{scaling}$ $\mathbf{head}$ processes the hidden states from the base model by applying fully connected linear layers to predict scalings. A softmax function is subsequently applied along the final dimension to ensure that the scalings across all experts sum to one. ReLU activations and dropout layers are interleaved between the linear layers within the X-LoRA scaling head. The hidden layer size is configurable and can be set to reduce the dimension from hidden size to scaling dimension effectively or adopt a widening-narrowing scheme akin to the feed-forward layer in transformer blocks.

The scalings form a matrix $\Lambda \in \mathbb{R}^{s \times l \times n}$, where $s$, $l$, and $n$ represent the sequence length, the number of layers, and the number of adapters, respectively. For each layer, the associated scalings are extracted as a matrix $\lambda \in \mathbb{R}^{s \times n}$. These scalings $\lambda$ are then applied to each LoRA adapter by broadcast multiplying the relevant X-LoRA scaling $\lambda_i \in \mathbb{R}^{1\times 1 \times s}$ with the input to the decomposition matrices $A$ and $B$.

Given $W_0 \in \mathbb{R}^{d \times k}$, $B_i \in \mathbb{R}^{d \times r_i}$, and $A_i \in \mathbb{R}^{r_i \times d}$, where the rank $r_i \ll {\rm min} (d,k)$ is unique to each adapter $i$, the forward pass of an X-LoRA layer is formulated as:

\begin{equation}
    h = W_0x + \sum_{i=1}^{n} B_iA_i(x \cdot  \lambda_i) \alpha_i
\end{equation}

Depictions of the scalings are presented in the main manuscript. These visualizations utilize either bar plots or the \texttt{contourf} function from Matplotlib~\cite{MatplotlibDocumentation}.

\subsection{X-LoRA implementation algorithm and code} The implementation selected here prioritizes ease of use and the flexibility to adapt the method to various models, especially those within the Hugging Face ecosystem. The X-LoRA algorithm functions by:

\begin{enumerate}
    \item Keeping the base model and adapters fixed, followed by tuning performance.
    \item Combining LoRA adapters based on predictions made by the X-LoRA scaling head.
\end{enumerate}

All code is developed in Python using PyTorch~\cite{Paszke2019PyTorchLibrary}. 

\subsubsection{Dual forward pass strategy for self-aware inference} Since X-LoRA needs to compute hidden states in order to predict the scaling factors, we adopt a dual forward pass approach. The first pass allows the model to assess the task and determine how to adjust itself, which is then followed by the actual inference step. This approach introduces a form of `self-awareness', enabling X-LoRA to allocate additional computation time for reflection rather than immediate response.

We refer to the first pass as the scaling pass and the second as the forward pass. During the scaling pass, the scalings $\Lambda$ are set to zero. We also tested setting the scalings $\Lambda$ to ${n}^{-1}$ and observed comparable outcomes, demonstrating stable performance (in the X-LoRA package, this parameter is configurable). To facilitate applying X-LoRA across a wide range of models, we implement a forward pass hook. This mechanism is crucial to coordinate the forward and scaling passes and is the key factor that allows compatibility with all models.

\begin{enumerate}
    \item $\mathbf{Scaling}$ $\mathbf{pass}$: Performed by running the base model with adapters fixed to constant values. The hidden states obtained are then used by the X-LoRA scaling head to produce $\Lambda$.
    \item $\mathbf{Forward}$ $\mathbf{pass}$: The scalings $\Lambda$ are applied to blend the adapters during the model’s forward pass. The resulting logits are returned as the final output of the X-LoRA model.
\end{enumerate}

As a caution, using the same key-value cache for both the scaling pass and forward pass may disrupt correct inputs to the scaling head, since the cache persists and accumulates across the two passes. Therefore, we avoid employing the key-value cache during training or inference, and we implemented proper management of this aspect in \texttt{Mistral.rs}.

\subsubsection{Freezing weights and enabling model segments for training} The base model and adapters remain frozen (i.e., their weights are not updated during training), making the X-LoRA scaling head the sole trainable component within an X-LoRA model. The code also provides the option to unfreeze all adapters along with the X-LoRA scaling head. Additionally, it is possible to train the entire model simultaneously—including the X-LoRA scaling head, adapters, and the foundational base model—potentially using different learning rates. Another avenue for future exploration could involve developing scaling functions that operate between two or more foundation models to realize a traceable implementation of SLERP-type modeling. This represents a promising direction for subsequent research.

\subsubsection{X-LoRA layers} An X-LoRA layer is designed to scale and then integrate the outputs from multiple adapters. Each X-LoRA layer encapsulates an original LoRA layer and thus can access the particular adapters associated with that layer.

To apply X-LoRA effectively to any model, the X-LoRA approach replaces the original LoRA layers by constructing X-LoRA layers and embedding their forward pass operations into these LoRA layers. This process utilizes features of Python’s object system to facilitate layer replacement.

In brief, during the forward pass, each X-LoRA layer performs the following operations (illustrated in Figure \ref{fig:general_arch}):

\begin{enumerate}
    \item $\mathbf{Scaling}$: Multiply the input signal of each LoRA adapter by its scaling coefficient $\lambda_i$.
    \item $\mathbf{Forward}$: Execute the forward pass through the wrapped LoRA layer.
\end{enumerate}

\subsection{Efficient Inference with \texttt{Mistral.rs} implemented in Rust} To enhance inference speed, we created \texttt{Mistral.rs}, a LLM serving platform written in Rust~\cite{matsakis2014rust} that supports X-LoRA and incorporates various performance optimizations. Specifically, the following techniques are implemented, presented here in no particular sequence:

\begin{itemize}
    \item LoRA adapter weight stacking: When all $A_i$ and $B_i$ matrices for each LoRA adapter share the same dimensions, it is possible to stack the LoRA adapter weight matrices $A_i$ and $B_i$ along their first dimension.
\end{itemize}

Mathematically, given $W_0 \in \mathbb{R}^{d \times k}$, for adapter matrices $A_i \in \mathbb{R}^{r \times d}$ and $B_i \in \mathbb{R}^{d \times r}$ with $r_i \ll {\rm min} (d,k)$, we stack the matrices so that $A \in \mathbb{R}^{i \times r \times d}$ and $B \in \mathbb{R}^{i \times d \times r}$. The $\alpha_i$ parameters are also applied to these stacked matrices during initialization.

To implement scalings, we rearrange the dimensions so that $\lambda \in \mathbb{R}^{i \times 1 \times 1}$, and then apply them via broadcast multiplication.

\item Non-granular scalings To minimize the frequency of running the scaling pass, we provide an option for non-granular scalings. This approach caches the scalings after $n$ completed runs, as the KV cache guarantees that the sequence length will remain one for the rest of the request. This method significantly reduces latency by lowering the number of calls to the base model.

\item Fused Compute Unified Device Architecture (CUDA) kernels To enhance the efficiency of the forward pass, we employ fused CUDA kernels. These fused kernels allow complex operations, such as Rotary Embedding (\cite{su2021roformer}), which are typically executed sequentially, to be merged into a single kernel capable of parallel execution on a GPU. Our kernels are based on the vLLM implementation (\cite{kwon2023efficient}).

\item Quantization The inference engine supports quantization, enabling the use of quantized base models. Initial experiments indicate that quantization up to 8 bits performs well even when using a model originally trained with \texttt{bfloat16}, for example.

\end{itemize}

\subsection{Datasets}

Each adapter is trained on specialized datasets, as summarized in Table~\ref{tab:table_loradata}, which also provides references to sources and original literature.

\subsection{Training strategy} Training proceeds in multiple stages. We start with the Zephyr-7B-$\beta$ model~\cite{HuggingFaceH4/zephyr-7b-betaFace}, which is built upon the Mistral-7B model \cite{Jiang2023Mistral7B}, and train a series of adapters using the datasets described above. The Zephyr-7B-$\beta$ tokenizer is used.

The first eight adapters are trained using question-answer pairs directly from the datasets (see Table~\ref{tab:table_loradata}).

The protein mechanics tasks follow a similar training approach but employ specific forward and inverse instruction sets. The bidirectional instruction tasks include:

\begin{quote}
\tiny {\texttt{CalculateForce< G E E C D C G S P ....> [0.310]\\
CalculateEnergy< G E E C D C G S P ....> [0.220]\\ 
CalculateForceHistory< G E E C D C G S P ....> [0.004,0.034,0.125,0.142,0.159,0.102,0.079,0.073,0.131, ...]\\
GenerateForce<0.310>\,[ G E E C D C G S P ....]\\
GenerateEnergy<0.220>\,[ G E E C D C G S P ....]\\
GenerateForceHistory<0.004,0.034,0.125,0.142,0.159,0.102,0.079,0.073,0.131, ...>\,[ G E E C D C G S P ....] }}
\end{quote}

These correspond to three distinct tasks in total: determining the peak unfolding force of a protein, calculating the unfolding energy, and recording the force-deformation profile. For each of these forward tasks, there are also inverse counterparts that produce a protein sequence as output, resulting in six tasks overall.

Each LoRA adapter is assigned a rank of $r=16$ with a scaling factor $\alpha=16$. Training is performed on five specific target modules: \texttt{v\_proj}, \texttt{k\_proj}, \texttt{q\_proj}, \texttt{gate\_proj}, and \texttt{o\_proj}, employing a dropout rate of 0.05. Typically, LoRA adapters are trained with relatively low ranks, ranging from 4 to 32, because higher ranks do not generally enhance performance but do increase computational load. Lower ranks yield more compact models with fewer parameters~\cite{hu2021lora}. The selected values for $r$, $\alpha$, and related settings are grounded in commonly adopted standards in the literature that have demonstrated efficacy. Additionally, prior investigations by the author of this paper tested various configurations and determined that the parameters chosen here represent a sensible compromise~\cite{Luu2023BioinspiredLLM:Materials}.

The X-LoRA scaling head is trained using approximately 2,000 samples, which are a subset extracted from the original dataset employed for training the adapters (additionally, we supplemented the training set with GPT-4 generated multiple choice question samples). We utilize a \texttt{paged\_adamw\_8bit} optimizer with a gradient clipping threshold of 0.3, a learning rate of $2\times 10^{-4}$ incorporating a warmup phase, and four steps of gradient accumulation, with a batch size set to one. The X-LoRA model undergoes training for roughly 10,000 steps, concluding with a loss of about 0.28.

Given that the base Mistral architecture applied here as the foundational model contains 32 layers, and considering that adapters are integrated into 5 layers per transformer layer, this results in a total of 160 LoRA layers, each replaced by an X-LoRA layer. The X-LoRA scaling head generates scaling factors for each of these 160 LoRA layers (for nine LoRA experts, this amounts to $160\times 9=1440$ $\lambda_i$ values computed per token). The model employs a single linear layer combined with ReLU activation functions and a dropout rate of $p=0.2$, encompassing roughly 5 million parameters. The implementation supports flexible architectures, allowing for the development and testing of more sophisticated deep classifier heads in various applications.

Both LoRA adapters and the X-LoRA scaling head are trained by predicting the next token using a cross-entropy loss function (we estimate the likelihood that the model assigns to the correct next token in a sequence based on preceding tokens; the training objective is to maximize this probability for the correct next token specified in the training data). We employ token shifting, where the sequence is offset by one token to generate the target sequence for the model’s prediction. For instance, if the input sequence is ``Proteins are fundamental to," the target training sequence would be ``are fundamental to biology." Thus, the model attempts to forecast the subsequent token in the sequence.

As depicted in Figure~\ref{fig:hierarchical_concept}, during training of the X-LoRA scaling head, only the scaling head parameters are updated while all other parameters remain fixed. The standard Zephyr tokenizer and prompt template are utilized.

\subsection{X-LoRA-Gemma model} The X-LoRA-Gemma model is constructed in a manner analogous to the X-LoRA model described above, but it is based on the Gemma-7B-it model~\cite{Google/gemma-7b-itFace}, employing four adapters (details in the main text). These four adapters are trained on datasets related to mechanics and materials, protein mechanics, bioinspired materials, as well as an additional quantum-mechanics-based molecular properties dataset QM9 (see Table~\ref{tab:table_QM9})~\cite{Ruddigkeit2012EnumerationGDB-17,Ramakrishnan2015ElectronicSpace}.

Each LoRA adapter has a rank of $r=16$ with $\alpha=16$. The training targets seven modules: \texttt{q\_proj}, \texttt{o\_proj}, \texttt{k\_proj}, \texttt{v\_proj}, \texttt{gate\_proj}, \texttt{up\_proj}, and \texttt{down\_proj}, applying a dropout rate of 0.05. The rank for each adapter is set to $r=16$ with $\alpha=16$. Given four adapters and seven target modules, and considering the 28 layers in the Gemma-7B base model, a total of $28\times 28=784$ $\lambda_i$ values are computed per token.

The first two adapters (bioinspired materials and mechanics of materials) incorporated in this model are trained using question-answer pairs provided by their respective datasets (see Table~\ref{tab:table_loradata}).

The protein mechanics adapters are trained similarly to the method outlined in the previous section, employing bidirectional forward and inverse instruction sets. Furthermore, we train an adapter using the QM9 dataset, focusing on two tasks:

\begin{quote}
\tiny {\texttt{CalculateMolecularProperties<CC(C)N1CC1C=O> [0.098,0.358,0.581,0.395,0.309,0.330,0.570,0.649,0.519,0.519,0.519,0.518]\\
GenerateMolecularProperties<0.098,0.358,0.581,0.395,0.309,0.330,0.570,0.649,0.519,0.519,0.519,0.518> [CC(C)N1CC1C=O]\\ 
...
}}
\end{quote}

These encompass two tasks in total: calculating molecular properties (the forward task) and the inverse task of generating a molecule with specified target properties. All 12 properties included in the QM9 dataset are either predicted or designed for, respectively (refer to Table~\ref{tab:table_QM9}).  
The model incorporates a total of four adapters.  

The X-LoRA scaling head of this model consists of three layers, with the intermediate layer sized at $3072 \times {FF}_{\rm fac} = 12,288$, where ${FF}_{\rm fac} = 4$ denotes a feed-forward multiplier. The architecture employs linear layers activated by ReLU functions and incorporates dropout with a probability of $p=0.2$, totaling approximately 47 million parameters.  

Training of the X-LoRA-Gemma scaling head utilized roughly 17,000 samples, which were sampled as a subset from the training dataset used for adapter training. Notably, although the X-LoRA-Gemma model contained only four adapters, the samples were drawn from the entire training set of the previously described model, supplemented with data from the QM9 dataset. This approach was taken to explore whether the X-LoRA model could acquire additional abilities during the scaling head training phase. Given that the loss converged well to 0.58 by training completion and demonstrated strong performance even on tasks outside its explicit training scope (e.g., chemistry tasks exemplified by the molecular design case discussed in the paper), we consider this method to generally produce favorable outcomes.

We employed a \texttt{paged\_adamw\_8bit} optimizer with a gradient clipping threshold of 0.3, a learning rate set at $1\times 10^{-4}$ including warmup, and utilized four gradient accumulation steps, with a batch size of one. The X-LoRA-Gemma model was trained for approximately 62,000 steps.

The original Gemma tokenizer along with its prompt template was used.

\subsection{Adversarial agentic modeling} We implemented an adversarial agentic approach by creating two X-LoRA agents. One agent specializes in question generation. It poses the initial question and continuously formulates further questions aimed at uncovering challenges and weaknesses within the responses. The other agent acts as a respondent, providing answers continuously. This approach was implemented using \{Guidance\}~\cite{Guidance-ai/guidance:Models.}.

Regarding the example connecting music and mechanics presented in the text, the characteristics of the two agents are described as follows. Firstly, the physicist and questioner:

\begin{quote}
\small\texttt{You are a critical physicist engaged in a discussion from a physics standpoint. Keep your replies concise and always challenge statements provocatively.

As a creative thinker, you incorporate ideas from other disciplines and push conceptual limits.}
\end{quote}

Secondly, a philosopher (in some instances this role is modified to represent a biology expert):

\begin{quote}
\small\texttt{You are a philosopher with expertise in biology, chemistry, and mathematics. You are participating in a discussion.

Keep your responses brief, yet accurate and inventive. You generate excellent ideas and new thought directions, maintaining logical consistency.}
\end{quote}

The questioner is specifically instructed to consider the current question and previous answers and to produce new, probing questions. This is implemented through a prompt structured as follows:

\begin{quote}
\small\texttt{Examine this question and answer.

\#\#\# Question: <question>

\#\#\# Response: <response>

\#\#\# Instruction: Provide a SINGLE follow-up question that critically examines the response.

Do NOT provide an answer or commentary at this stage.

The single question is: }
\end{quote}

After the conversation has reached $N$ turns, the model is required to 1) summarize the main insights discussed, 2) enumerate the key insights as bullet points, and 3) highlight the most significant takeaway.

The raw conversation text serves as the foundation for further analysis through graph construction.

\subsection{Knowledge graph generation} We employ Zephyr-7B-$\beta$ to extract triplets from the text, following the methodology described in~\cite{Rahulnyk/knowledge_graph:QnA} with enhancements based on the Llama Index graph generation method~\cite{Liu_LlamaIndex_2022}. Graphs are displayed using NetworX\cite{Networkx/networkx:Python} and Pyvis~\cite{WestHealth/pyvis:Graphs.}. A central instruction provided aligns with the approach in~\cite{Liu_LlamaIndex_2022}:

\begin{quote}
\small\texttt{Your objective is to extract the ontology of terms referenced in the given context. These terms should capture the principal concepts according to the context.

Present your output as a JSON list. Each list entry should include a pair of terms and the relationship between them, for example: [...] }
\end{quote}

We guide the construction of triplets consisting of nodes and edge attributes as follows, adhering to the method from \cite{Liu_LlamaIndex_2022}:

\begin{quote}
\small\texttt{\begin{itemize}
            \item "node\_1": "A concept from the extracted ontology"
            \item "node\_2": "A related concept from the extracted ontology"
            \item "edge": "A concise description of the relationship between node\_1 and node\_2"
        \end{itemize}
        }
\end{quote}

We also supply the model with an example of ontological analysis, such as:

\begin{quote}
\small\texttt{
       Context: ```Spider silk is a strong natural fiber used to catch prey in a web.``` }
\end{quote}

Subsequently, example triplets are provided, for instance:

\begin{quote}
\small\texttt{\begin{itemize}
            \item "node\_1": "spider silk"
            \item "node\_2": "fiber"
            \item "edge": "is"
        \end{itemize}
        }
\end{quote}

The assembled graph data is then grouped into communities using the Girvan-Newman algorithm~\cite{Girvan2002CommunityNetworks}.

\subsection{Visualization of molecular structures}

We utilize PyMOL~\cite{Schrodinger2015The1.8} to visualize and examine the predicted protein structures. Various built-in scripts and functions within the software enable identification of specific protein features, such as secondary structure, hydrophobic/hydrophilic areas, disulfide bridges, and hydrogen bonding.

\section*{Data Availability Statement} The codes and datasets that underpin the results of this work are publicly accessible at \url{https://github.com/EricLBuehler/xlora} and \url{https://huggingface.co/lamm-mit/x-lora}. Additional data sources referenced in this study are cited (see especially Table~\ref{tab:table_loradata}). 

\end{document}